\documentclass[10pt,a4paper]{amsart}
\usepackage[margin=19mm]{geometry}
\usepackage{amsmath,amssymb,mathtools}
\usepackage[scr=rsfs,cal=euler]{mathalfa}
\usepackage{xfrac}
\usepackage[unicode,hidelinks,bookmarks=false]{hyperref}
\newcommand{\ie}{i.e.\ }
\newcommand{\cf}{cf.\ }
\newcommand{\resp}{respectively\ }
\newcommand{\Subsection}{\subsection}
\providecommand{\nobreakdash}{\nobreak\hbox{-}}
\setlength{\emergencystretch}{2em}
\multlinegap0pt
\usepackage{bm}
\usepackage{bbm}
\usepackage{dsfont}
\usepackage{xspace}
\usepackage{cancel}
\usepackage{mathtools}
\usepackage{amsthm}
\makeatletter
\@ifundefined{theorem}{\newtheorem{theorem}{Theorem}}{}
\makeatother
\newcommand{\op}[1]{\operatorname{#1}}
\title[On the average $2$-torsion in class groups and narrow class ]{On the average $2$-torsion in class groups and narrow class groups of cubic orders with prescribed shape}
\author{Anwesh Ray}
\date{Source: arXiv:2601.04503v1 (CC BY 4.0)}
\begin{document}
\maketitle

\noindent\emph{Source and licence.} On the average $2$-torsion in class groups and narrow class groups of cubic orders with prescribed shape. Version arXiv:2601.04503v1 (CC BY 4.0), distributed under the Creative Commons Attribution 4.0 International licence, as recorded in the arXiv licence metadata for this exact version and shown on the abstract page https://arxiv.org/abs/2601.04503v1. Full rights notice: licenses/arxiv-2601.04503-CC-BY-4.0.txt.\par\smallskip

\noindent\textit{Editorial scope.} Counted unit: the Theorem whose source label is \texttt{None} in the article (source0.tex lines 389--394, source label \texttt{None}), delivered together with the article's own complete original proof of it (source0.tex lines 396--416, one \texttt{proof} environment of 21 lines). No mathematics is written, repaired or re-derived: statement and proof are the article's own text line for line. Required context retained uncounted: none. Every definition, display and label of those lines is retained unchanged. Omitted from this package and not redistributed: 1--388, 395--395, 417--843. Nothing inside a retained excerpt is omitted or altered: every retained body is delivered exactly as the article's lines are written, so no omission receipt of any kind accompanies this package. Citations: the retained excerpts carry no citation command at all, so no bibliography entry of the article is transcribed and no reference is redistributed. Reference adaptation: none was needed and none was made; every reference inside the delivered unit resolves inside it, so no unresolved reference or citation is produced. Printed slips kept literal: no printed word, symbol, hypothesis or number of the counted statement or of its proof was altered. Layout and class: the article's own class is \texttt{amsart} with packages including geometry, amscd, amsmath, amsxtra, amsthm, amssymb, stmaryrd, xr, mathrsfs, mathtools, enumerate, commath, comment, orcidlink; the wrapper is the lane's pinned \texttt{amsart} editorial wrapper at 10pt on A4, which restates the theorem-like environments the retained text needs and adds the article's own macro definitions that the retained text needs (\texttt{\textbackslash op}), with extra wrapper packages (bm, bbm, dsfont, xspace, cancel, mathtools). The delivered numbering is the unit's own. Assets: no retained span contains a figure environment, a graphics-inclusion command, a PDF-page inclusion, a table or a TikZ picture; no image file is copied into this package, the package declares no asset dependency and no image file is redistributed with it. Licence: version arXiv:2601.04503v1 of this article is distributed under the Creative Commons Attribution 4.0 International licence, as recorded in the arXiv licence metadata for this exact version and shown on the abstract page https://arxiv.org/abs/2601.04503v1. A full rights notice is delivered at licenses/arxiv-2601.04503-CC-BY-4.0.txt. Characters: every retained excerpt is pure ASCII, so the delivered engine, XeLaTeX with its default Latin Modern text font, reports no missing character.
\begin{theorem}[Delone--Faddeev, Gan--Gross--Savin]
There is a canonical bijection between isomorphism classes of cubic rings $R$ over $T$ and
$\op{GL}_{2}(T)$-equivalence classes of binary cubic forms $f$ with coefficients in $T$.
Under this correspondence, let $R(f)$ be the cubic ring associated to a form $f$. The discriminant of $R(f)$ coincides with the
discriminant of $f$. In particular, \(R(f)\) is nondegenerate if and only if \(\op{Disc}(f)\neq 0\).
\end{theorem}

\begin{proof} We provide a sketch of the proof for the benefit of exposition. To describe the correspondence explicitly, let \(R\) be a cubic ring. Choosing a \(T\)-basis \(\langle 1,\omega,\theta\rangle\) for \(R\), and translating \(\omega\) and \(\theta\) by suitable integers, we may assume that \(\omega\theta\in T\). Such a basis is called \emph{normal}. With respect to a normal basis, multiplication in \(R\) is given by relations of the form
\[
\omega\theta = n, \qquad
\omega^{2} = m + b\omega - a\theta, \qquad
\theta^{2} = \ell + d\omega - c\theta,
\]
for integers \(a,b,c,d,\ell,m,n\). Associativity of multiplication imposes constraints on \(\ell,m,n\), and a direct calculation shows that these constants are uniquely determined by
\[
n=-ad, \qquad m=-ac, \qquad \ell=-bd.
\]
The coefficients \(a,b,c,d\) therefore determine the cubic ring \(R\) uniquely, and one associates to \(R\) the binary cubic form
\[
f(x,y)=ax^{3}+bx^{2}y+cxy^{2}+dy^{3}.
\]

Conversely, starting from a binary cubic form \(f(x,y)=ax^{3}+bx^{2}y+cxy^{2}+dy^{3}\), the above relations define a cubic ring \(R(f)\) with underlying \(T\)-module \(T\oplus T\omega\oplus T\theta\). The associativity relations ensure that this construction is well-defined and inverse to the assignment \(R\mapsto f\). A direct computation shows that the discriminant of the cubic ring \(R(f)\) coincides with the discriminant of the associated binary cubic form,
\[
\op{Disc}(R(f))=\op{Disc}(f)
= b^{2}c^{2}-4ac^{3}-4b^{3}d-27a^{2}d^{2}+18abcd.
\]
\end{proof}

\end{document}
